\documentclass[12pt]{article}
\usepackage{amsmath, amssymb}
\usepackage{physics}
\usepackage{xcolor}
\usepackage[margin=1in]{geometry}

\title{From Wave Equation to Helmholtz Equation: \\
Complete Derivation for Inhomogeneous Media}
\author{Detailed Explanation}
\date{\today}

\begin{document}

\maketitle

\section{The Question}

\textbf{How do we get the Helmholtz equation from the wave equation in inhomogeneous media?}

This is a crucial step that is often glossed over in textbooks. Let's work through it carefully.

\section{Starting Point: Wave Equation in Inhomogeneous Media}

From Maxwell's equations in an inhomogeneous, source-free dielectric medium, we derived:
\begin{equation}
\nabla^2 \mathbf{E} - \mu_0 \epsilon(r) \frac{\partial^2 \mathbf{E}}{\partial t^2} = \mu_0 \frac{\partial \Phi}{\partial t} \nabla \epsilon(r)
\end{equation}

For simplicity, let's first consider the case where $\nabla \epsilon(r) = 0$ locally (i.e., in regions away from interfaces):
\begin{equation}
\nabla^2 \mathbf{E} - \mu_0 \epsilon(r) \frac{\partial^2 \mathbf{E}}{\partial t^2} = 0
\end{equation}

Using $\epsilon(r) = \epsilon_0 n^2(r)$ and $c = 1/\sqrt{\mu_0 \epsilon_0}$:
\begin{equation}
\boxed{\nabla^2 \mathbf{E} - \frac{n^2(r)}{c^2} \frac{\partial^2 \mathbf{E}}{\partial t^2} = 0}
\end{equation}

This is the \textbf{wave equation in an inhomogeneous medium}.

\section{Key Assumption: Monochromatic Fields}

\subsection{Time-Harmonic Ansatz}

We now make the critical assumption that the field is \textbf{monochromatic} (single frequency):
\begin{equation}
\mathbf{E}(r, t) = \mathbf{E}(r) e^{-i\omega t}
\end{equation}

\textbf{Why is this valid?}
\begin{itemize}
\item Any field can be decomposed into monochromatic components via Fourier transform
\item Linear systems: solve for each frequency separately, then superpose
\item For optical fibers, laser sources are approximately monochromatic
\end{itemize}

\subsection{Substituting into Wave Equation}

Let's substitute $\mathbf{E}(r, t) = \mathbf{E}(r) e^{-i\omega t}$ into the wave equation.

\textbf{Step 1: First time derivative}
\begin{equation}
\frac{\partial \mathbf{E}}{\partial t} = \frac{\partial}{\partial t}\left[\mathbf{E}(r) e^{-i\omega t}\right] = -i\omega \mathbf{E}(r) e^{-i\omega t}
\end{equation}

\textbf{Step 2: Second time derivative}
\begin{equation}
\frac{\partial^2 \mathbf{E}}{\partial t^2} = \frac{\partial}{\partial t}\left[-i\omega \mathbf{E}(r) e^{-i\omega t}\right] = (-i\omega)^2 \mathbf{E}(r) e^{-i\omega t} = -\omega^2 \mathbf{E}(r) e^{-i\omega t}
\end{equation}

\textbf{Step 3: Spatial derivatives}

The spatial derivatives only act on $\mathbf{E}(r)$, not on $e^{-i\omega t}$:
\begin{equation}
\nabla^2 \mathbf{E}(r,t) = \nabla^2 \left[\mathbf{E}(r) e^{-i\omega t}\right] = \left[\nabla^2 \mathbf{E}(r)\right] e^{-i\omega t}
\end{equation}

\subsection{The Derivation}

Substituting into Eq. (3):
\begin{equation}
\left[\nabla^2 \mathbf{E}(r)\right] e^{-i\omega t} - \frac{n^2(r)}{c^2} \left[-\omega^2 \mathbf{E}(r) e^{-i\omega t}\right] = 0
\end{equation}

Factor out $e^{-i\omega t}$:
\begin{equation}
e^{-i\omega t} \left[\nabla^2 \mathbf{E}(r) + \frac{\omega^2 n^2(r)}{c^2} \mathbf{E}(r)\right] = 0
\end{equation}

Since $e^{-i\omega t} \neq 0$, we can divide both sides by it:
\begin{equation}
\nabla^2 \mathbf{E}(r) + \frac{\omega^2 n^2(r)}{c^2} \mathbf{E}(r) = 0
\end{equation}

\section{Final Form: Helmholtz Equation}

Define the \textbf{free-space wavenumber}:
\begin{equation}
k_0 = \frac{\omega}{c} = \frac{2\pi}{\lambda_0}
\end{equation}

Then:
\begin{equation}
\frac{\omega^2}{c^2} = k_0^2
\end{equation}

Substituting:
\begin{equation}
\boxed{\nabla^2 \mathbf{E}(r) + k_0^2 n^2(r) \mathbf{E}(r) = 0}
\end{equation}

This is the \textbf{Helmholtz equation for inhomogeneous media}.

\section{Component Form}

For each Cartesian component:
\begin{equation}
\nabla^2 E_i(r) + k_0^2 n^2(r) E_i(r) = 0, \quad i = x, y, z
\end{equation}

\section{Physical Interpretation}

\subsection{Comparison with Homogeneous Case}

\textbf{Homogeneous medium} ($n(r) = n$ = constant):
\begin{equation}
\nabla^2 \mathbf{E} + k^2 \mathbf{E} = 0, \quad k = k_0 n
\end{equation}
Solutions: plane waves $\mathbf{E} = \mathbf{E}_0 e^{i\mathbf{k} \cdot \mathbf{r}}$

\textbf{Inhomogeneous medium} ($n(r)$ varies with position):
\begin{equation}
\nabla^2 \mathbf{E} + k_0^2 n^2(r) \mathbf{E} = 0
\end{equation}
Solutions: more complex (modes, guided waves, etc.)

\subsection{The Role of $n^2(r)$}

The term $n^2(r)$ acts as a \textbf{position-dependent eigenvalue}:
\begin{itemize}
\item In the core ($r < a$): $n^2(r) = n_{\text{core}}^2$ (larger)
\item In the cladding ($r > a$): $n^2(r) = n_{\text{clad}}^2$ (smaller)
\end{itemize}

This creates a \textbf{potential well} for the electromagnetic field, leading to guided modes.

\section{Special Case: Step-Index Fiber}

For a step-index fiber:
\begin{equation}
n^2(r) = \begin{cases}
n_{\text{core}}^2, & r < a \\
n_{\text{clad}}^2, & r > a
\end{cases}
\end{equation}

The Helmholtz equation becomes:
\begin{equation}
\nabla^2 \mathbf{E} + k_0^2 n^2(r) \mathbf{E} = 0
\end{equation}

with different eigenvalues in core vs. cladding:
\begin{itemize}
\item \textbf{Core}: $\nabla^2 \mathbf{E} + k_0^2 n_{\text{core}}^2 \mathbf{E} = 0$
\item \textbf{Cladding}: $\nabla^2 \mathbf{E} + k_0^2 n_{\text{clad}}^2 \mathbf{E} = 0$
\end{itemize}

\section{Solving the Helmholtz Equation in Cylindrical Coordinates}

\subsection{Cylindrical Laplacian}

In cylindrical coordinates $(r, \phi, z)$:
\begin{equation}
\nabla^2 = \frac{1}{r}\frac{\partial}{\partial r}\left(r \frac{\partial}{\partial r}\right) + \frac{1}{r^2}\frac{\partial^2}{\partial \phi^2} + \frac{\partial^2}{\partial z^2}
\end{equation}

\subsection{Separation of Variables}

For fiber modes, we assume:
\begin{equation}
\mathbf{E}(r, \phi, z) = \mathbf{E}(r, \phi) e^{i\beta z}
\end{equation}

where $\beta$ is the propagation constant along the fiber axis.

Substituting into Helmholtz equation:
\begin{equation}
\frac{1}{r}\frac{\partial}{\partial r}\left(r \frac{\partial E}{\partial r}\right) + \frac{1}{r^2}\frac{\partial^2 E}{\partial \phi^2} + (i\beta)^2 E + k_0^2 n^2(r) E = 0
\end{equation}

Simplifying:
\begin{equation}
\frac{1}{r}\frac{\partial}{\partial r}\left(r \frac{\partial E}{\partial r}\right) + \frac{1}{r^2}\frac{\partial^2 E}{\partial \phi^2} + \left[k_0^2 n^2(r) - \beta^2\right] E = 0
\end{equation}

For azimuthally symmetric modes ($\partial/\partial \phi = 0$):
\begin{equation}
\frac{1}{r}\frac{d}{dr}\left(r \frac{dE}{dr}\right) + \left[k_0^2 n^2(r) - \beta^2\right] E = 0
\end{equation}

\section{Summary of the Derivation Chain}

\begin{center}
\begin{tabular}{|c|l|}
\hline
\textbf{Step} & \textbf{Equation} \\
\hline
1. Maxwell's equations & $\nabla \times \mathbf{E} = -\partial \mathbf{B}/\partial t$ \\
& $\nabla \times \mathbf{H} = \partial \mathbf{D}/\partial t$ \\
\hline
2. Wave equation & $\nabla^2 \mathbf{E} - \frac{n^2(r)}{c^2} \frac{\partial^2 \mathbf{E}}{\partial t^2} = 0$ \\
\hline
3. Time-harmonic ansatz & $\mathbf{E}(r,t) = \mathbf{E}(r) e^{-i\omega t}$ \\
\hline
4. Helmholtz equation & $\nabla^2 \mathbf{E}(r) + k_0^2 n^2(r) \mathbf{E}(r) = 0$ \\
\hline
5. Cylindrical coordinates & Separation of variables in $(r, \phi, z)$ \\
\hline
6. Eigenvalue equation & Boundary conditions $\Rightarrow$ discrete $\beta_m$ \\
\hline
\end{tabular}
\end{center}

\section{Key Points to Remember}

\begin{enumerate}
\item The transformation $\partial^2/\partial t^2 \to -\omega^2$ comes from the time-harmonic ansatz

\item The $n^2(r)$ factor remains position-dependent in the Helmholtz equation

\item In the fiber core and cladding separately, $n^2$ is constant, giving different Bessel-type solutions

\item Boundary conditions at $r = a$ determine the allowed propagation constants $\beta_m$

\item The Helmholtz equation is an \textbf{eigenvalue problem}: given $\omega$ (or $k_0$), find $\beta$ and $\mathbf{E}(r)$
\end{enumerate}

\section{Why This Matters for Quantum Communication}

The Helmholtz equation modal solutions are essential because:

\begin{itemize}
\item They define the \textbf{mode functions} $f_m(r)$ used in field quantization:
\begin{equation}
\hat{\mathbf{A}}(r) = \sum_m \sqrt{\frac{\hbar \omega_m}{2\epsilon_0 \epsilon_m V_m}} \left[\hat{a}_m f_m(r) + \hat{a}_m^\dagger f_m^*(r)\right]
\end{equation}

\item The propagation constant $\beta_m$ determines:
\begin{itemize}
\item Phase velocity: $v_p = \omega/\beta$
\item Group velocity: $v_g = d\omega/d\beta$
\item Dispersion: $\beta_2 = d^2\beta/d\omega^2$
\end{itemize}

\item The effective mode area determines the coupling strength:
\begin{equation}
g \propto \frac{1}{\sqrt{V_{\text{eff}}}} \propto \frac{1}{a}
\end{equation}
\end{itemize}

\end{document}